% !TEX program = pdflatex
% Scanned source: single_amplitude_formula/pictures/ct_f_a_a_scan.png
% The double commutator of the CT expansion: 1/2[[F,A],A] = 1/2[C,T] + h.c.,
% and its expansion into the eight commutators of the sheets.
\documentclass[11pt]{article}
\usepackage[T1]{fontenc}
\usepackage[utf8]{inputenc}
\usepackage{lmodern}
\usepackage{amsmath,amssymb}
\usepackage[margin=1in]{geometry}
\usepackage{booktabs}
\usepackage[hidelinks]{hyperref}

\newcommand{\E}[2]{\hat E^{#1}_{#2}}
\newcommand{\hc}{\mathrm{h.c.}}
\newcommand{\occ}{i}
\newcommand{\vir}{a}
\newcommand{\Cd}{\hat C}

\title{\bfseries The double commutator of the CT expansion\\[2pt]
       \large $\tfrac12[[\hat F,\hat A],\hat A]$, and the eight blocks}
\author{}
\date{}

\begin{document}
\maketitle
\vspace{-2.5em}

\noindent
This is the last piece of the first-order structure: the second-order term of
the CT expansion, in the convention of \texttt{ct\_F\_A.tex} and
\texttt{ct\_V\_A.tex} (crossed Fock order, excitations, general indices over
$\{\occ\}\cup\{\vir\}$). With
\begin{equation}
  \hat A=\hat T-\hat T^\dagger,
  \qquad
  \hat T=\hat T_1+\hat T_2,
  \qquad
  \hat C=[\hat F,\hat A],
  \label{eq:setup}
\end{equation}
the second-order term is
\begin{equation}
  \tfrac12\bigl[[\hat F,\hat A],\hat A\bigr]
   =\tfrac12[\hat C,\hat A]
   =\tfrac12[\hat C,\hat T]+\tfrac12[\hat C,\hat T]^\dagger,
  \label{eq:first}
\end{equation}
the second equality because $\hat A$ is anti-Hermitian, which makes $\hat C$
Hermitian, so that $[\hat C,\hat T]^\dagger=-[\hat C,\hat T^\dagger]$ and
$[\hat C,\hat A]=[\hat C,\hat T]-[\hat C,\hat T^\dagger]$. Only the $\hat T$
half has to be computed; the other half is its conjugate.

\section*{1. Splitting \texorpdfstring{$\hat C$}{C}}

Write the two halves of the Fock commutator separately,
\begin{equation}
  \hat C_1=[\hat F,\hat T_1],
  \qquad
  \hat C_2=[\hat F,\hat T_2],
  \qquad
  \hat C=\hat C_1+\hat C_2+\hat C_1^\dagger+\hat C_2^\dagger,
  \label{eq:Csplit}
\end{equation}
with $\hat C_r^\dagger=-[\hat F,\hat T_r^\dagger]$ the conjugate of
$\hat C_r$. Both $\hat C_1$ and $\hat C_2$ are already in closed form in
\texttt{ct\_F\_A.tex}: two terms for the one-body block, four for the two-body
block.

Since the commutator is linear in its first argument,
\begin{equation}
  [\hat C,\hat T]
   =[\hat C_1,\hat T]+[\hat C_2,\hat T]
    +[\hat C_1^\dagger,\hat T]+[\hat C_2^\dagger,\hat T],
  \label{eq:four}
\end{equation}
and expanding each $[\cdot,\hat T]=[\cdot,\hat T_1]+[\cdot,\hat T_2]$ gives the
eight commutators of the sheets:
\begin{equation}
  \boxed{\;
  \begin{aligned}
  \tfrac12\bigl[[\hat F,\hat A],\hat A\bigr]
   =\tfrac12\Bigl\{&
    [\hat C_1,\hat T_1]+[\hat C_1,\hat T_2]
   +[\hat C_2,\hat T_1]+[\hat C_2,\hat T_2]\\
   &+[\hat C_1^\dagger,\hat T_1]+[\hat C_1^\dagger,\hat T_2]
   +[\hat C_2^\dagger,\hat T_1]+[\hat C_2^\dagger,\hat T_2]\Bigr\}
   +\hc .
  \end{aligned}\;}
  \label{eq:eight}
\end{equation}

\section*{2. What the eight blocks are}

The first four of \eqref{eq:eight} are the nested double commutators of the
Fock operator with two $\hat T$'s,
\begin{equation}
  [\hat C_r,\hat T_s]=\bigl[[\hat F,\hat T_r],\hat T_s\bigr],
  \qquad r,s\in\{1,2\},
  \label{eq:nested}
\end{equation}
and the last four are the same with the conjugate inner commutator,
$[\hat C_r^\dagger,\hat T_s]$. They are not related to each other by
conjugation in a way that saves work:
\begin{equation}
  [\hat C_r,\hat T_s]^\dagger=-[\hat C_r^\dagger,\hat T_s^\dagger],
  \label{eq:adjoint}
\end{equation}
which is a statement about $[\cdot,\hat T_s^\dagger]$, not about
$[\cdot,\hat T_s]$, so each of the eight has to be evaluated. In particular the
two mixed nested commutators are different operators,
\begin{equation}
  \bigl[[\hat F,\hat T_1],\hat T_2\bigr]\neq
  \bigl[[\hat F,\hat T_2],\hat T_1\bigr],
\end{equation}
since $\hat T_1\hat T_2-\hat T_2\hat T_1$ does not vanish (four terms), so the
Jacobi shortcut is not available and both mixed blocks are kept.

Each block is evaluated the same way: commute once to get $\hat C_r$ (or
$\hat C_r^\dagger$) from \texttt{ct\_F\_A.tex}, then commute the result with
$\hat T_s$ using the $(1,1)$ or $(1,2)$ rule. Counting terms before merging:
\begin{center}
\footnotesize
\begin{tabular}{@{}llrc@{}}
\toprule
block & inner object & inner terms & terms out \\
\midrule
$[\hat C_1,\hat T_1]=\bigl[[\hat F,\hat T_1],\hat T_1\bigr]$
  & $\hat C_1$ (1 body) & 2 & $2\times2=4$ \\
$[\hat C_1,\hat T_2]=\bigl[[\hat F,\hat T_1],\hat T_2\bigr]$
  & $\hat C_1$ (1 body) & 2 & $2\times4=8$ \\
$[\hat C_2,\hat T_1]=\bigl[[\hat F,\hat T_2],\hat T_1\bigr]$
  & $\hat C_2$ (2 body) & 4 & $4\times2=8$ \\
$[\hat C_2,\hat T_2]=\bigl[[\hat F,\hat T_2],\hat T_2\bigr]$
  & $\hat C_2$ (2 body) & 4 & $4\times4=16$ \\
the four conjugate-inner blocks & $\hat C_r^\dagger$ & same & same \\
\bottomrule
\end{tabular}
\end{center}

\noindent
Every one of these terms is kept: as in the first-order blocks, the terms of a
block are different operators with different index patterns, and the collection
into $\bar h$ and $\bar g$ reads a coefficient off each of them.

\subsection*{2.1 A worked block: $\bigl[[\hat F,\hat T_1],\hat T_1\bigr]$}

The simplest of the eight, and the one that answers the natural question of
whether it vanishes. With $[\hat F,\hat T_1]$ from \texttt{ct\_F\_A.tex} and
$\hat T_1=t^b_j\E{b}{j}$ with $b$ virtual and $j$ occupied,
\begin{equation}
  \bigl[[\hat F,\hat T_1],\hat T_1\bigr]
   =f^\vir_q\,t^\vir_\occ\,t^\beta_j
      \bigl[\E{q}{\occ},\E{\beta}{j}\bigr]
    -f^p_\occ\,t^\vir_\occ\,t^\beta_j
      \bigl[\E{\vir}{p},\E{\beta}{j}\bigr],
  \label{eq:worked}
\end{equation}
and the $(1,1)$ rule gives
\begin{equation}
  \bigl[\E{q}{\occ},\E{\beta}{j}\bigr]
   =\delta^\beta_\occ\E{q}{j}-\delta^q_j\E{\beta}{\occ},
  \qquad
  \bigl[\E{\vir}{p},\E{\beta}{j}\bigr]
   =\delta^\beta_p\E{\vir}{j}-\delta^\vir_j\E{\beta}{p}.
  \label{eq:workedrule}
\end{equation}
Two of the four resulting terms vanish identically, and not by cancellation:
$\delta^\beta_\occ$ ties an occupied label to a virtual one, and so does
$\delta^\vir_j$, so both are zero whatever the amplitudes are. The other two
survive and are equal to each other under $\occ\leftrightarrow j$,
$\vir\leftrightarrow\beta$,
\begin{equation}
  \bigl[[\hat F,\hat T_1],\hat T_1\bigr]
   =-f^\vir_j\,t^\vir_\occ\,t^\beta_j\,\E{\beta}{\occ}
    -f^\beta_\occ\,t^\vir_\occ\,t^\beta_j\,\E{\vir}{j}
   =-2\,f^\vir_j\,t^\vir_\occ\,t^\beta_j\,\E{\beta}{\occ}.
  \label{eq:workedresult}
\end{equation}
So the block does \emph{not} vanish. $[\hat F,\hat T_1]$ is a one-body operator
of the same shape as $\hat T_1$ but with different coefficients, and two
one-body operators do not commute merely because they look alike. Checked on
matrices: the four delta terms of \eqref{eq:workedrule} reproduce the block to
$2.2\times10^{-15}$, and the block is non-zero.

\section*{3. The check}

The identities of \S1 are operator identities, so they can be checked on
explicit matrices with no reference to the $\hat E$ algebra.
\path{tests/test_ct_structure.py} does this with random Hermitian $\hat F$ and
arbitrary $\hat T_1,\hat T_2$; the values are
\begin{center}
\footnotesize
\begin{tabular}{@{}lr@{}}
\toprule
identity & max$|\Delta|$ \\
\midrule
$\hat C=[\hat F,\hat A]=\hat C_1+\hat C_2+\hat C_1^\dagger+\hat C_2^\dagger$
  & $3.7\times10^{-14}$ \\
$[\hat C,\hat T]$ equals the four commutators of \eqref{eq:four}
  & $6.8\times10^{-13}$ \\
$[\hat C,\hat T]$ equals the eight of \eqref{eq:eight}
  & $7.6\times10^{-13}$ \\
$\tfrac12[[\hat F,\hat A],\hat A]=\tfrac12([\hat C,\hat T]+\hc)$
  & $4.8\times10^{-13}$ \\
$[\hat C_1,\hat T_1]=\bigl[[\hat F,\hat T_1],\hat T_1\bigr]$
  & $0.0$ \\
$[\hat C_1,\hat T_2]^\dagger=-[\hat C_1^\dagger,\hat T_2^\dagger]$
  & $1.7\times10^{-13}$ \\
\bottomrule
\end{tabular}
\end{center}
and the check also asserts that
$\bigl[[\hat F,\hat T_1],\hat T_2\bigr]$ and
$\bigl[[\hat F,\hat T_2],\hat T_1\bigr]$ differ, so the mixed blocks cannot be
taken as equal.

\section*{4. Where this leaves the derivation}

With \eqref{eq:eight} the second-order term is reduced to the eight blocks
above, each of which is two applications of rules that are already verified.
What remains is not algebra of this kind but two other things:
\begin{itemize}
  \item \textbf{the three-body reduction.} The $(2,2)$-derived blocks produce
        three-body operators, which the $[\dots]_{1,2}$ truncation folds back
        into one- and two-body form with the extended normal ordering of
        Kutzelnigg and Mukherjee --- an approximation, and one that has to be
        recorded as such wherever the result is used;
  \item \textbf{the collection.} Reading $\bar h$ and $\bar g$ off the
        assembled monomials, term by term, in the general basis
        $C^q_p\E{q}{p}+\tfrac12C^{rs}_{pq}\E{rs}{pq}$ of \texttt{ct\_F\_A.tex},
        which is where the index patterns of the previous notes have to be
        preserved rather than cancelled.
\end{itemize}

\begin{thebibliography}{9}
\bibitem{Kutzelnigg:1997} W.~Kutzelnigg and D.~Mukherjee,
  \emph{Normal order and extended Wick theorem for a multiconfiguration
  reference wave function}, J.~Chem.~Phys.\ \textbf{107}, 432 (1997).
\bibitem{Yanai:2012cj} T.~Yanai and T.~Shiozaki,
  \emph{Canonical transcorrelated theory with projected Slater-type geminals},
  J.~Chem.~Phys.\ \textbf{136}, 084107 (2012).
\end{thebibliography}

\end{document}
